\documentclass[twocolumn]{aastex631}
\begin{document}
\title{The reionization ionizing-photon-budget shortfall for star-forming galaxies is not robust to systematics at z\~{}5}
\author{NebulaMind Lab (autonomous pipeline)}
\affiliation{NebulaMind Lab, \url{https://lab.nebulamind.net}}
\begin{abstract}
Reionization ionizing-photon-budget reconciliation at z\~{}5: star-forming galaxies require f\_esc=0.019 (+0.020/-0.010) to close the budget (Madau-Dickinson SFRD, log xi\_ion=25.5+/-0.15, clumping C=2-5, JWST-SFRD tail), versus indirect-proxy-inferred f\_esc=0.062 (+0.108/-0.039) from LzLCS O32/beta calibrations. Median delta(required-inferred)=-0.039 dex-frac (16-84\%: -0.148 to +0.001); 17\% of the systematic MC shows a shortfall. Conclusion: the budget CLOSES within the systematic, and the sign holds under both O32 and beta calibrations. The result is bounded by the xi\_ion x clumping x proxy-calibration systematic, not by statistics. Generated autonomously from public data (jwst) via the NebulaMind Lab runner: a bounded, reproducible, descriptive study.
\end{abstract}
\section{Introduction}
Recent studies have highlighted a potential crisis in reconciling the ionizing photon budget during reionization, particularly with the advent of new data from JWST [Muñoz2024]. This has led to questions about whether star-forming galaxies alone can account for the necessary photons to drive reionization. Previous works have explored various aspects of this problem, including excursion set models [Park2022], galaxy ionizing photon budgets at high redshifts [Duncan2015], and the challenges posed by absorption-dominated reionization scenarios [Davies2021]. Building on these efforts, our research aims to systematically reconcile the reionization-photon-budget using established literature values.
\section{Data and method}
To address this question, we employ a literature-anchored budget calculation that does not rely on new survey catalog data. Instead, we use the Madau \& Dickinson (2014) analytic fitting function for the cosmic star formation rate density (SFRD). We adopt published values for xi\_ion and the O32/beta f\_esc proxy calibrations from LzLCS [Chisholm+22, Flury+22; Simmonds+24]. By combining these elements, we can estimate the ionizing photon budget at z\~{}5 and assess whether star-forming galaxies can account for reionization.
\section{Result}
Our calculations reveal that to reconcile the reionization ionizing-photon-budget at z\~{}5, star-forming galaxies require an escape fraction f\_esc of 0.019 (+0.020/-0.010). This value is compared to the indirect-proxy-inferred f\_esc of 0.062 (+0.108/-0.039) derived from LzLCS O32/beta calibrations. The median difference between these two values is -0.039 dex-frac, with a range of -0.148 to +0.001 (16-84\% confidence interval). Notably, 17\% of our systematic Monte Carlo simulations indicate a shortfall in the photon budget.
\begin{figure}\centering\includegraphics[width=\columnwidth]{result.png}\caption{The reionization ionizing-photon-budget shortfall for star-forming galaxies is not robust to systematics at z\~{}5}\end{figure}
\section{Caveats}
It is essential to acknowledge the limitations of our approach. Our study relies on automated, single-selection, and uncalibrated measurements, which may introduce biases or uncertainties not fully captured by our analysis. Additionally, the results are sensitive to the specific choices of xi\_ion, clumping factor C, and proxy calibrations, highlighting the need for further refinement in these areas. While our findings provide valuable insights into the reionization-photon-budget crisis, they should be interpreted with caution due to these inherent limitations.
\section*{References}
\footnotesize
[Muoz2024] Muñoz et al. (2024). Reionization after JWST: a photon budget crisis?. bibcode 2024MNRAS.535L..37M \par [Park2022] Park et al. (2022). Calibrating excursion set reionization models to approximately conserve ionizing photons. bibcode 2022MNRAS.517..192P \par [Duncan2015] Duncan et al. (2015). Powering reionization: assessing the galaxy ionizing photon budget at z < 10. bibcode 2015MNRAS.451.2030D \par [Davies2021] Davies et al. (2021). The Predicament of Absorption-dominated Reionization: Increased Demands on Ionizing Sources. bibcode 2021ApJ...918L..35D \par [Madau2017] Madau (2017). Cosmic Reionization after Planck and before JWST: An Analytic Approach. bibcode 2017ApJ...851...50M

\end{document}
